\section{Odd-$k$ Separation for All Odd $k$}
\label{sec:odd-k-separation}

The two-point $1$-NN witness shows the separation for $k=1$.  A natural
question is whether the same phenomenon persists for larger odd $k$.  The
answer is yes: a general construction works for every odd $k \ge 3$.

\subsection{Construction}

Let $k = 2m+1$ for an integer $m \ge 1$.  Let the metric space consist of two
distinct points $a$ and $b$ with $d(a,b) > 0$.  Define the ordered sample
\[
S_m \;=\; \big((a,0)_1,\; \ldots,\; (a,0)_{m+1},\;
          (a,1)_1,\; \ldots,\; (a,1)_{m},\;
          (b,0)_1,\; \ldots,\; (b,0)_{m}\big),
\]
where subscripts track the occurrence index within the ordered tuple.  The
sample has size $n = 3m+1$.

Fix the query-label pair $(x,y) = (a,0)$ and the neighborhood size $k = 2m+1$.

\subsection{Original margin}

At query $a$, the ordered neighbors are sorted by distance then occurrence
index.  All occurrences at $a$ are at distance $0$; all occurrences at $b$ are
at distance $d(a,b) > 0$.  The deterministic order therefore selects all
$(m+1)$ copies of $(a,0)$ first, then all $m$ copies of $(a,1)$, giving exactly
$k = 2m+1$ neighbors.  The signed margin is
\[
\M_k(S_m, a) = (m+1)(-1) + m(+1) = -1 < 0,
\]
so the classifier predicts $0$ at $a$.  Hence the loss at $(a,0)$ is $0$.

\subsection{LOO case analysis}

We show $\Delta_{\mathrm{loo}}(S_m, i) = 0$ for every index $i$, covering the
three types of occurrence.

\paragraph{LOO at an $(a,0)$ occurrence.}
Delete the $i$-th copy of $(a,0)$ and evaluate at the deleted occurrence
$(a,0)$.  After deletion the remaining sample has size $n-1 = 3m$, and
$k' = \min(k, n-1) = 2m+1$ for all $m \ge 1$.

At query $a$, the top-$k'$ neighbors consist of the remaining $m$ copies of
$(a,0)$, all $m$ copies of $(a,1)$, and the first copy of $(b,0)$.  The margin
is
\[
\M_{k'}(S_m^{-i}, a) = m(-1) + m(+1) + 1(-1) = -1.
\]
The prediction is $0$, unchanged from the original.  Since the loss at $(a,0)$
was $0$ before and remains $0$ after deletion, $\Delta_{\mathrm{loo}} = 0$.

\paragraph{LOO at an $(a,1)$ occurrence.}
Delete the $i$-th copy of $(a,1)$ and evaluate at $(a,1)$.  After deletion, the
top-$k'$ at $a$ consists of all $(m+1)$ copies of $(a,0)$, the remaining
$(m-1)$ copies of $(a,1)$, and the first copy of $(b,0)$.  The margin is
\[
\M_{k'}(S_m^{-i}, a) = (m+1)(-1) + (m-1)(+1) + 1(-1) = -3.
\]
The prediction is $0$.  The original prediction at $(a,1)$ was also $0$ (margin
$-1$), so the original loss at $(a,1)$ was $1$; the loss after deletion is
still $1$.  Hence $\Delta_{\mathrm{loo}} = 0$.

\paragraph{LOO at a $(b,0)$ occurrence.}
Delete the $i$-th copy of $(b,0)$ and evaluate at $(b,0)$.  The original margin
at query $b$ is
\[
\M_k(S_m, b) = m(-1) + (m+1)(-1) = -(2m+1),
\]
so the original prediction is $0$ and the loss at $(b,0)$ is $0$.

After deletion, the top-$k'$ at $b$ consists of $(m-1)$ copies of $(b,0)$, all
$(m+1)$ copies of $(a,0)$, and the first copy of $(a,1)$.  The margin is
\[
\M_{k'}(S_m^{-i}, b) = (m-1)(-1) + (m+1)(-1) + 1(+1) = -(2m-1).
\]
For all $m \ge 0$, this is $\le 0$, so the prediction remains $0$ and the loss
stays $0$.  Thus $\Delta_{\mathrm{loo}} = 0$.

Since every index falls into one of the three cases above, we have
$\Delta_{\mathrm{loo}}(S_m, i) = 0$ for all $i$.

\subsection{Replacement-induced margin crossing}

Now replace the first occurrence $(a,0)_1$ with $(a,1)$.  The perturbed sample
is
\[
S_m^{1 \leftarrow (a,1)}
= \big((a,1),\; (a,0)_2,\ldots,(a,0)_{m+1},\;
        (a,1)_1,\ldots,(a,1)_m,\;
        (b,0)_1,\ldots,(b,0)_m\big).
\]

At query $a$, the top-$k$ neighbors are now: $(a,1)$ at index $1$, then
$(a,0)_2$ through $(a,0)_{m+1}$ ($m$ copies), then $(a,1)_1$ through
$(a,1)_m$ ($m$ copies).  All $k = 2m+1$ neighbors are at distance $0$.  The
signed margin becomes
\[
\M_k(S_m^{1 \leftarrow (a,1)}, a) = 1(+1) + m(-1) + m(+1) = 1 > 0,
\]
so the classifier now predicts $1$ at $a$.  The loss at $(a,0)$ changes from
$0$ to $1$, giving $\Delta_{\mathrm{rep}}(S_m, 1, (a,1), a, 0) = 1$.

\begin{theorem}[Odd-$k$ LOO/replace-one separation]
\label{thm:odd-k-separation}
For every odd $k = 2m+1$ with $m \ge 1$, the sample $S_m$ satisfies
$\Delta_{\mathrm{loo}}(S_m, i) = 0$ for every index $i$, while there exists an
index $j$ and a replacement point $z$ such that $\Delta_{\mathrm{rep}}(S_m, j,
z, a, 0) = 1$.
\end{theorem}

The $k=1$ case, excluded from \Cref{thm:odd-k-separation}, is already covered
by the two-point witness of \Cref{sec:minimal-1nn}.  Taken together, the
LOO/replace-one separation holds for every odd $k$: $k=1$ by the two-point
witness and $k \ge 3$ by \Cref{thm:odd-k-separation}.

\subsection{Interpretation}

The construction uses only the fact that $a$ and $b$ are distinct points with
$d(a,b) > 0$; the separation therefore holds in any metric space with at least
two points.  The ordered occurrence sequence is essential: placing all $(a,0)$
copies before $(a,1)$ copies ensures the original margin is $-1$.  Conflicting
labels at $a$ are the mechanism by which a single replacement changes the label
composition of the top-$k$ neighborhood without changing its geometric
membership.

The construction stabilizes every LOO evaluation point using duplicate
occurrences, but keeps the designated query exactly one vote away from a
replacement-induced margin crossing.  This is the unifying pattern behind all
the separation results in this paper.
